% !Mode:: "TeX:UTF-8"
% !TEX program  = xelatex
\section{系统模型与问题定义}
\subsection{系统模型}
本文考虑一个多租户的大语言模型推理系统，其中一个共享的基础模型被加载至 GPU 上，而多个 LoRA 适配器用于支持不同任务或用户请求。设适配器集合为 $\mathcal{A}=\{a_1,a_2,\dots,a_N\}$，系统在运行过程中需要处理一系列请求序列 $\mathcal{R}=\{r_1,r_2,\dots\}$，其中每个请求 $r_i$ 关联唯一的适配器 $a(r_i)\in\mathcal{A}$。

由于 GPU 显存资源有限，系统无法同时容纳所有适配器。设 GPU 可用于适配器的显存预算为 $M$，每个适配器 $a$ 占用显存大小为 $\mathrm{mem}(a)$，则任意时刻驻留在 GPU 上的适配器集合 $\mathcal{H}_t\subseteq\mathcal{A}$ 必须满足如下约束：
\begin{equation}
\sum_{a\in\mathcal{H}_t}\mathrm{mem}(a)\le M
\end{equation}

当一个请求到达时，若其所需适配器已经存在于当前驻留集合 $\mathcal{H}_t$ 中，则可以直接执行推理；否则，系统需要将该适配器从 CPU 或外部存储加载至 GPU。该过程通常伴随着显著的时间开销，并可能触发已有适配器的淘汰操作。此外，不同适配器在显存占用与加载开销上可能存在显著差异，这进一步增加了资源管理的复杂性。

\subsection{问题定义}
在上述系统模型下，本文关注如下核心问题：在给定请求序列 $\mathcal{R}$ 和显存约束 $M$ 的条件下，如何动态管理适配器的驻留与淘汰，以优化系统整体性能。具体而言，在每个时间步 $t$，系统接收到一个请求 $r_t$，其对应适配器为 $a_t=a(r_t)$。系统需要根据当前状态（包括历史请求、当前驻留集合以及系统负载情况）决定是否需要加载新的适配器，并在必要时选择合适的适配器进行淘汰。
\begin{equation}
\min\sum_t\mathbb{I}_{\text{miss}}(t)\cdot\mathrm{cost}(a_t)
\end{equation}

\subsection{基线方法}
\subsubsection{FIFO (First-In, First-Out)}
FIFO 按照适配器进入 GPU 的时间顺序进行管理。显存不足时，优先淘汰最早加载的适配器。其实现简单，但无法利用访问模式信息。

\subsubsection{LRU (Least Recently Used)}
LRU 假设近期访问的适配器未来更可能被再次使用。系统维护最近访问时间戳并在显存不足时淘汰“最久未使用”者。

\subsubsection{LFU (Least Frequently Used)}
LFU 基于累计访问频率进行淘汰。该策略能较好识别长期热点，但缺乏时间敏感性，在热点迁移场景中适应较慢。

\subsubsection{方法对比分析}
上述方法均属于被动缓存策略，仅基于历史访问信息进行决策，未显式建模未来请求，也未充分考虑适配器显存占用差异。

\section{方法设计}
\subsection{方法概述}
本文提出一种基于预测的适配器调度方法。在显存受限条件下，系统通过对适配器未来使用趋势进行轻量级预测，动态维护高价值适配器驻留集合。

\subsection{基于 EMA 的适配器使用预测}
对于任意适配器 $a$，其在时间步 $t$ 的 EMA 值定义为：
\begin{equation}
EMA_t(a)=\lambda\cdot EMA_{t-1}(a)+(1-\lambda)\cdot\mathbb{I}(a_t=a)
\end{equation}
其中 $\lambda\in[0,1]$ 为衰减因子，$\mathbb{I}(a_t=a)$ 为指示函数。

\subsection{适配器效用建模}
\begin{equation}
S_t(a)=\alpha\cdot EMA_t(a)+\beta\cdot Q_t(a)+\gamma\cdot A_t(a)
\end{equation}
其中 $Q_t(a)$ 表示队列中等待该适配器的请求数，$A_t(a)$ 为活跃状态指示量，$\alpha,\beta,\gamma$ 为权重参数。

\subsection{显存感知的效用归一化}
\begin{equation}
\tilde{S}_t(a)=\frac{S_t(a)}{(\mathrm{mem}(a))^\delta}
\end{equation}
归一化使系统偏向“单位显存收益更高”的适配器。

\subsection{动态适配器选择策略}
\begin{equation}
\mathcal{H}_t^*=\underset{\mathcal{H}\subseteq\mathcal{A}}{\operatorname{argmax}}\sum_{a\in\mathcal{H}}\tilde{S}_t(a),\quad \text{s.t.}\sum_{a\in\mathcal{H}}\mathrm{mem}(a)\le M
\end{equation}
该问题类似 0-1 背包，在线场景下采用基于效用排序的贪心近似求解。

\subsection{在线决策过程}
\textbf{Algorithm 1:} 基于预测的适配器调度算法。
\begin{enumerate}
  \item 更新请求对应适配器的 EMA 及统计信息；
  \item 计算所有适配器效用评分 $\tilde{S}_t(a)$；
  \item 若请求适配器已驻留，则直接执行；
  \item 否则在必要时按低效用优先淘汰，直至满足显存约束后加载目标适配器；
  \item 更新驻留集合并进入下一步请求。
\end{enumerate}
